\section*{Hoofdstuk \ref{ch:b1:proteins} --- Aminozuren en eiwitten}

\begin{solution}{exo:b1:proteins:1}
$\mathrm{^+H_3N{-}CH(R){-}COO^-}$: een \omterm{def:b1:proteins:aminoacid}{zwitterion}. Leu apolair alifatisch; Ser
\omterm{def:b1:water-small-molecules:water}{polair} ongeladen; Lys positief geladen; Glu negatief geladen; Phe aromatisch.
\end{solution}

\begin{solution}{exo:b1:proteins:2}
Primair: de sequentie, \omterm{def:b1:proteins:peptide}{peptidebindingen}. Secundair: de plaatselijke vouwing van de
ruggengraat (helix, plaat), \omterm{def:b1:water-small-molecules:water}{waterstofbruggen} van de ruggengraat. Tertiair: de
vouwing van de hele keten; \omterm{prop:b1:water-small-molecules:properties}{hydrofoob effect}, \omterm{def:b1:water-small-molecules:water}{waterstofbruggen}, zoutbruggen,
vanderwaalscontacten, \omterm{def:b1:proteins:tertiary}{disulfidebruggen}. Quaternair: de samenvoeging van
subeenheden, door dezelfde niet-covalente krachten.
\end{solution}

\begin{solution}{exo:b1:proteins:3}
$450\times 110 = \qty{50}{kDa}$. Helix: $450\times 0.15 = \qty{68}{nm}$.
Streng: $450\times 0.35 = \qty{158}{nm}$.
\end{solution}

\begin{solution}{exo:b1:proteins:4}
Ongeveer $0.98$ en $0.76$: een vijfde van de lading (\qty{22}{\%}) wordt
afgeleverd.
\end{solution}

\begin{solution}{exo:b1:proteins:5}
Groepen: de N-terminale aminogroep (9.5), de \omterm{def:b1:proteins:aminoacid}{zijketen} van Lys (10.5), de \omterm{def:b1:proteins:aminoacid}{zijketen}
van Asp (3.9), de C-terminale carboxylgroep (2). pH 7: $+1 + 1 - 1 - 1 = 0$.
pH 2: $+1 + 1 + 0 - 0.5 = +1.5$ (de eindstandige carboxylgroep half geïoniseerd).
pH 12: $0 + 0 - 1 - 1 = -2$ (Lys grotendeels gedeprotoneerd). De netto lading is
nul tussen de p$K_a$ van Asp (3.9) en die van de aminogroep (9.5), waar de lading
overal $0$ is: p$I \approx 6.7$ (het gemiddelde van 3.9 en 9.5).
\end{solution}

\begin{solution}{exo:b1:proteins:6}
De \omterm{def:b1:proteins:aminoacid}{zijketen} van proline is aan zijn eigen stikstof terug gebonden: de stikstof
draagt geen waterstof om de \omterm{def:b1:water-small-molecules:water}{waterstofbrug} van de helix te geven, en $\phi$ ligt
vast op ongeveer $-60^\circ$, wat alleen in de eerste winding van een helix wordt
verdragen. Glycine, met een waterstof als \omterm{def:b1:proteins:aminoacid}{zijketen}, ondervindt geen botsing en kan
waarden van $\phi$ en $\psi$ aannemen die elk ander residu verboden zijn, ook die
van scherpe bochten.
\end{solution}

\begin{solution}{exo:b1:proteins:7}
Paringen van acht cysteïnen: $7\times 5\times 3\times 1 = 105$; willekeurige
heroxidatie geeft de natieve reeks in ongeveer \qty{1}{\%} van de gevallen. Dat de
activiteit vrijwel volledig terugkeerde, toonde aan dat de keten en niet het toeval
de paring koos: de sequentie legt de vouwing vast en de disulfiden zetten haar
alleen maar vast.
\end{solution}

\begin{solution}{exo:b1:proteins:8}
$P_{50}^{2.8} = 33.4$. Bij 2: $2^{2.8} = 6.96$, $Y = 0.17$; bij 3.5: $0.50$; bij
8: $8^{2.8} = 337$, $Y = 0.91$. Met $n = 1$ geldt $Y =
P/(3.5 + P)$: $0.36$, $0.50$, $0.70$. Aflevering tussen 8 en 2: \omterm{def:b1:proteins:allostery}{coöperatief}
$0.74$, \omterm{def:b1:proteins:allostery}{niet-coöperatief} $0.34$: het coöperatieve \omterm{def:b1:proteins:peptide}{eiwit} levert twee keer zoveel af.
\end{solution}

\begin{solution}{exo:b1:proteins:9}
Een tetrameer van vier identieke subeenheden van \qty{50}{kDa}. Laat de SDS-gel
zonder reducerend middel lopen: verschijnen er banden bij 100, 150 of
\qty{200}{kDa}, dan verbinden \omterm{def:b1:proteins:tertiary}{disulfidebruggen} de subeenheden; blijft alleen de
band van \qty{50}{kDa} over, dan worden zij niet-covalent bijeengehouden.
\end{solution}

\begin{solution}{exo:b1:proteins:10}
Glutamaat is geladen en gehydrateerd aan het oppervlak; valine is hydrofoob en
schept een kleverige plek die water uitsluit. In de gedeoxygeneerde T-toestand ligt
op een naburig molecuul een complementaire hydrofobe holte bloot, zodat de
moleculen kop aan staart tot vezels aaneenplakken die de rode \omterm{def:b1:cell-unit-of-life:cell}{cel} vervormen. In de
longen verbergt de R-toestand de holte en lossen de vezels op; in de \omterm{def:b1:body-plans-tissues:tissue}{weefsels}, waar
hemoglobine gedeoxygeneerd is, ontstaan zij --- daarom sikkelen de \omterm{def:b1:cell-unit-of-life:cell}{cellen} daar en
verstoppen zij de kleine bloedvaten.
\end{solution}

\begin{solution}{exo:b1:proteins:11}
Bij \qty{37}{\celsius} is $RT = \qty{2.58}{kJ/mol}$: $K = e^{-40/2.58} =
e^{-15.5} = \num{1.8e-7}$: één molecuul op vijf miljoen is ontvouwen. Bij
\qty{45}{\celsius} met \qty{10}{kJ/mol} is $RT = 2.64$: $K = e^{-3.8} =
0.022$: \qty{2}{\%} ontvouwen, en dat loopt steil op. Een krappe stabiliteit laat
\omterm{def:b1:proteins:peptide}{eiwitten} van vorm veranderen wanneer zij werken (\omterm{def:b1:proteins:allostery}{allosterie}, katalyse), ontvouwd
worden voor transport door membranen, en snel worden afgebroken en vervangen
wanneer zij beschadigd of overbodig zijn.
\end{solution}

\begin{solution}{exo:b1:proteins:12}
Anfinsen toonde dat de sequentie volstaat om de vouwing vast te leggen, en de
\omterm{def:b1:proteins:allostery}{allosterie} toont dat de vorm --- en de verandering ervan --- de functie is. Maar de
\omterm{def:b1:cell-unit-of-life:cell}{cel} grijpt bij elke stap in: \omterm{prop:b1:proteins:denaturation}{chaperonnes} redden ketens die anders zouden klonteren
voordat zij gevouwen zijn, want het volgepakte \omterm{def:b1:eukaryotic-cell:organelle}{cytosol} is geen verdunde
reageerbuis; één enkele vervanging (sikkelcel) geeft een correct gevouwen \omterm{def:b1:proteins:peptide}{eiwit}
waarvan het oppervlak verkeerd is, zodat ``vorm'' de oppervlaktechemie moet
omvatten en niet alleen de ruggengraat; en veel \omterm{def:b1:proteins:peptide}{eiwitten} hebben na de vouwing nog
modificaties, partners of liganden nodig voordat zij werken. De zin verwoordt het
beginsel; de \omterm{def:b1:cell-unit-of-life:cell}{cel} levert de voorwaarden.
\end{solution}

\begin{solution}{pb:b1:proteins:1}
\textbf{1.} $150\times 1.34 = \qty{201}{mL}$ $\mathrm{O_2}$ per liter.
\textbf{2.} $e^{3.51} = 33.4$; $e^{4.67} = 106.7$; $e^{7.25} = 1408$.
\textbf{3.} Longen: $1408/(33.4 + 1408) = 0.977$. \omterm{def:b1:body-plans-tissues:tissue}{Weefsels}:
$106.7/(33.4 + 106.7) = 0.762$.
\textbf{4.} $0.977 - 0.762 = 0.215$: $0.215\times 201 = \qty{43}{mL}$ per liter.
\textbf{5.} $250/43 = \qty{5.8}{L/min}$ --- dicht bij het hartminuutvolume in
rust.
\textbf{6.} Myoglobine: longen $13.3/13.67 = 0.973$; \omterm{def:b1:body-plans-tissues:tissue}{weefsels}
$5.3/5.67 = 0.935$; het levert $0.038\times 201 = \qty{7.7}{mL}$ per liter af, zes
keer minder: een hyperbolische drager die goed laadt kan bij weefseldrukken niet
lossen.
\textbf{7.} \omterm{def:b1:proteins:allostery}{Coöperativiteit} legt het steile deel van de curve tussen de twee
werkdrukken, zodat een kleine drukdaling een groot deel van de lading vrijgeeft.
\textbf{8.} $e^{2.78} = 16.1$; $Y = 16.1/(33.4 + 16.1) = 0.325$.
\textbf{9.} $e^{4.21} = 67.4$; $Y = 16.1/(67.4 + 16.1) = 0.193$.
\textbf{10.} Zonder Bohr: $(0.977 - 0.325)\times 201 = \qty{131}{mL}$;
met: $(0.977 - 0.193)\times 201 = \qty{158}{mL}$ per liter; winst \qty{27}{mL},
\qty{20}{\%}.
\textbf{11.} $2.7/3.07 = 0.88$ verzadigd; bij \qty{0.5}{kPa} is
$0.5/0.87 = 0.57$: het geeft tijdens de samentrekking een derde van zijn voorraad
aan de \omterm{def:b1:eukaryotic-cell:mitochondrion}{mitochondriën} af, en laadt tussen de samentrekkingen weer bij.
\textbf{12.} Protonen binden op plaatsen (histidinen, de N-uiteinden) weg van de
heemgroepen, en stabiliseren door dat binden de T-toestand, waardoor de affiniteit
bij de heemgroepen daalt: een ligand op de ene plaats verandert de affiniteit op
een andere --- \omterm{def:b1:proteins:allostery}{allosterie}.
\textbf{13.} $e^{5.45} = 233$; longen $233/(33.4 + 233) = 0.875$; \omterm{def:b1:body-plans-tissues:tissue}{weefsels}
$0.762$; aflevering $0.113\times 201 = \qty{23}{mL}$: ongeveer de helft van
zeeniveau.
\textbf{14.} $e^{3.88} = 48.4$: longen $233/(48.4 + 233) = 0.828$; \omterm{def:b1:body-plans-tissues:tissue}{weefsels}
$106.7/(48.4 + 106.7) = 0.688$; aflevering $0.140\times 201 =
\qty{28}{mL}$.
\textbf{15.} Het laden daalt een beetje (0.875 naar 0.828) maar het lossen daalt
meer (0.762 naar 0.688), zodat het verschil groeit: bij \qty{7}{kPa} liggen de
longen nog op de vlakke top van de curve terwijl de \omterm{def:b1:body-plans-tissues:tissue}{weefsels} op het steile deel
zitten.
\textbf{16.} Capaciteit $180\times 1.34 = \qty{241}{mL}$; $0.140\times
241 = \qty{34}{mL}$ per liter --- drie kwart van zeeniveau hersteld.
\textbf{17.} $6^{2.8} = e^{2.8\times 1.79} = e^{5.02} = 151$: longen
$233/384 = 0.607$, \omterm{def:b1:body-plans-tissues:tissue}{weefsels} $106.7/258 = 0.414$: aflevering $0.193$, op papier nog
beter --- maar het slagaderlijk bloed zou dan slechts \qty{61}{\%} verzadigd zijn,
en elke verdere daling van de longdruk of elke vraag om een hogere weefseldruk zou
de hersenen tekortdoen; de werkelijke aanpassing houdt op rond \qty{4}{kPa}.
\textbf{18.} $48.4/(33.4 + 48.4) = 0.59$.
\textbf{19.} $e^{2.68} = 14.6$; $48.4/(14.6 + 48.4) = 0.77$.
\textbf{20.} Bij dezelfde druk bindt het foetale \omterm{def:b1:proteins:peptide}{eiwit} meer: terwijl het bloed van
de moeder tot \qty{59}{\%} lost en het foetale bloed tot \qty{77}{\%} laadt,
stroomt zuurstof de drukgradiënt af die het verschil in affiniteit over het
placentamembraan in stand houdt.
\textbf{21.} BPG bindt de T-toestand en verlaagt de affiniteit; de
$\gamma$-ketens binden het zwak, zodat foetaal hemoglobine dichter bij zijn eigen
hoge affiniteit blijft: een lagere $P_{50}$.
\textbf{22.} Longen: $1408/(14.6 + 1408) = 0.990$; \omterm{def:b1:body-plans-tissues:tissue}{weefsels} $106.7/(14.6
+ 106.7) = 0.880$; aflevering $0.110\times 201 = \qty{22}{mL}$, de helft van die
van de volwassene: goed om zuurstof van de moeder over te nemen, slecht om haar af
te leveren zodra de longen werken, vandaar de omschakeling.
\textbf{23.} $5^{2.8} = e^{2.8\times 1.609} = e^{4.51} = 90.6$: longen
$1408/1499 = 0.94$, \omterm{def:b1:body-plans-tissues:tissue}{weefsels} $106.7/197 = 0.54$: aflevering $0.40$, bijna het
dubbele van normaal --- de persoon is licht blauw aan de lippen maar verdraagt dat
goed; de prijs is een kleinere reserve op hoogte.
\textbf{24.} $Y = P/(3.5 + P)$: longen $0.79$, \omterm{def:b1:body-plans-tissues:tissue}{weefsels} $0.60$: aflevering
$0.19\times 201 = \qty{38}{mL}$ --- en in de sprintende spier bij \qty{2.7}{kPa}:
$0.44$; aflevering $0.35\times 201 = \qty{70}{mL}$ tegen
\qty{158}{mL}. Zonder \omterm{def:b1:proteins:allostery}{coöperativiteit} laadt de drager slecht en lost hij slecht.
\textbf{25.} \qty{43}{mL} $\mathrm{O_2}$ per liter in rust, \qty{158}{mL} per
liter in een sprintende spier; het verschil komt van de \omterm{def:b1:proteins:allostery}{coöperativiteit} (de
sigmoïde curve) en van het bohreffect (de verschuiving van $P_{50}$ in \omterm{def:b1:water-small-molecules:ph}{zuur}).
\end{solution}
